\section{Experimental Methodology}
\label{sec:methodology}

The three concerns evaluated in this paper---context management
(\S\ref{sec:context}), credential isolation (\S\ref{sec:credentials}), and
production hardening (\S\ref{sec:hardening})---share one measurement basis.
We state it once, here, and each concern section cites it rather than
re-deriving it.

\subsection{Provider-authoritative measurement}
\label{subsec:method-proxy}

Every LLM call made by every framework under test is routed through a
single, shared LiteLLM proxy sitting in front of the model provider. Token
counts and dollar costs are read from the proxy's own accounting of the
request and response it forwarded, not from any figure a framework
self-reports. This choice removes an entire class of measurement error: a
framework's internal token estimate can be stale, provider-specific, or
simply wrong, but the proxy sees the literal request that left the process
and the literal response that came back, and it is that transcript we
measure.

Each proxy call is attributed to the run that issued it. For the five
Python competitor frameworks, which support custom request headers, we tag
every outbound call with a per-run identifier header and have the proxy
route matching spans to a run-scoped log. Colmena's OpenAI-compatible
client adapter does not expose a header-injection point, so its calls are
instead correlated to a run by a session-file line-count delta: the number
of new spans the proxy has recorded since the run began is read off the
shared span log rather than off a header the call itself carries. Both
methods attribute every call to exactly one run before any aggregate figure
is computed.

The proxy's span record is deliberately narrow: it captures token counts,
timestamps, and the model identifier for each call, and nothing of the
message content itself. The evidence this paper reports is therefore
always a token \emph{count}, never an inspection of transcript text; where
a later section states that "the token counts are the evidence"
(\S\ref{sec:context}), this is the reason.

\subsection{Frameworks and conditions}
\label{subsec:method-frameworks}

We evaluate six frameworks: Colmena, whose execution engine is written in
Rust (\S\ref{sec:model}), and five Python frameworks in wide production
use: CrewAI, LangChain, LangGraph, LlamaIndex, and Google ADK. All six are
held to the same experimental conditions. Each uses the same model alias,
\texttt{gemini-2.5-flash}, at temperature $0$, against the same proxy
endpoint, and, within each concern's task, the same inputs. No framework
is given an easier prompt, a different tool surface, or a different
model.

Within those fixed conditions, each framework is exercised
\emph{idiomatically}: we construct each agent using that framework's own
default conventions for memory, context retention, and tool invocation,
rather than a conservative or hand-tuned configuration common to all six.
This is a deliberate methodological choice, not an oversight. The paper's
claim is about what a framework does \emph{by default}, for an author who
follows its own documentation, and an idiomatic construction is the only
construction that tests that claim. Where a competitor's default can be
matched by hand-architected code, we say so explicitly and evaluate that
steelman as a separate, clearly labeled arm (\S\ref{subsec:context-closure},
\S\ref{subsec:cred-steelman})---never folded into the default comparison.

\subsection{Price is a held constant; the measured variable is volume}
\label{subsec:method-price}

Because every framework is driven through the same proxy against the same
model alias, the price charged per token is identical across all six
arms of every comparison in this paper. We hold it constant by
construction: no framework pays a different rate for an input or output
token than any other. Consequently, every dollar figure this paper reports
is a direct, fixed-rate function of a token count, and the quantity that
varies---and the quantity this paper measures---is token \emph{volume}:
how many tokens a framework's default construction sends and receives to
accomplish the same task. A cost table is a token-count table read through
a shared multiplier, not an independent measurement.

We note, for scope only, that a provider offering prompt caching could
lower the realized dollar cost of re-sending previously seen context
without changing the underlying token count; such an optimization is
orthogonal to what we measure and does not appear in any run reported
here. This is not a caveat about the validity of the comparison---price
per token is a constant by construction in every arm, not a variable that
could favor or disadvantage any framework, and it is not treated as a
limitation anywhere in this paper.

\subsection{Quality guardrail}
\label{subsec:method-guardrail}

A framework that spends fewer tokens by producing a worse answer has not
demonstrated an efficiency gain; every comparison in this paper is
therefore gated on task quality before any token or leakage figure is
taken to be meaningful. We fix a single guardrail, in advance of the runs,
and apply it identically to every framework and every concern. On each
task's designated scoring turns, the framework's answer is checked by an
automated substring test against a fixed set of ground-truth strings
established when the task was authored: the answer must contain the
expected substrings for that turn to count as a pass. This check is
mechanical and requires no judgment call at scoring time, which is why it
is the criterion of record. We report, alongside it, a complementary score
from an LLM judge that rates each answer's factual fidelity to the task's
source material on a continuous scale, as a secondary corroborating
signal. Each concern section reports its frameworks' pass or fail against
this criterion before reporting the token, cost, or capability comparison
that motivates the section.

\subsection{Reproducibility}
\label{subsec:method-repro}

All Colmena figures in this paper are produced by the engine build tagged
\texttt{colmena\_dag\_engine-v0.9.0}. Every run reported---for Colmena and
for the five competitor frameworks alike---is driven by a shared benchmark
harness: one proxy configuration, one span-logging mechanism, and one
per-experiment driver script per task, invoked identically for every
framework under test. A given experiment's task definition, scoring
logic, and driver script are therefore the same artifact whether the
framework being measured is Colmena or a competitor, which is what makes
the comparisons in \S\ref{sec:context}--\S\ref{sec:hardening} a controlled
measurement of the frameworks rather than of the harness that runs them.
